\subsection{Ablation: Competitor Diversity and Weighting}
\label{sec:competitor_ablation}

How should a limited discovery budget be distributed across competitors?
We separate three choices: weighting observed pairs, varying the number of
competitors, and selecting their semantic families. All variants retain
paired subtraction and the cosine/top-5\% readout.

\paragraph{Protocol.}
We use frozen Qwen2.5-VL-3B-Instruct and two saved discovery collections:
\emph{Core} (10 categories, 2,104 records, 45 observed pairs) and
\emph{Extended} (76 categories, 16,609 records, 304 observed pairs).
We refit directions from the saved contrasts and rescore all variants on
the same reconstructed evaluation images: 5,000 COCO val2017 images and,
for the additional OpenImages categories, 25,099 OpenImages images.
Unverified OpenImages image--category labels are excluded.
L18 is primary. Sampling comparisons average per-fit metrics over 20
seeds on a common eligible category cohort. Confidence intervals (CIs)
use 2,000 paired image-bootstrap draws within each source dataset and
condition on these fixed fits; seed SD is reported separately.
This is a follow-up analysis on previously examined validation data.

% Generated by scripts/render_competitor_ablation.py; do not edit numbers.
\begin{table}[t]
    \centering
    \footnotesize
    \setlength{\tabcolsep}{4pt}
    \begin{tabular}{lrrr}
        \toprule
        Menu (categories) & $K$ & Macro AUC $\uparrow$ & AP $\uparrow$ \\
        \midrule
        Core (7) & 2 & $0.9610\pm0.0083$ & 0.8168 \\
        Core (7) & 4 & $0.9668\pm0.0048$ & 0.8510 \\
        Core (7) & 8 & $\mathbf{0.9691}\pm0.0023$ & \textbf{0.8654} \\
        Extended (64) & 2 & $0.8426\pm0.0111$ & 0.6099 \\
        Extended (64) & 4 & $0.8524\pm0.0094$ & 0.6350 \\
        Extended (64) & 8 & $\mathbf{0.8564}\pm0.0069$ & \textbf{0.6432} \\
        \bottomrule
    \end{tabular}
    \caption{\textbf{Competitor count at a fixed budget.} Qwen2.5-VL-3B-Instruct, L18, 32 fitting records per category. AUC is mean $\pm$ sample SD over 20 fitting seeds; AP is the mean macro AP. All values use the 0--1 scale. Each menu uses a common eligible cohort across $K$. Seed SD measures fitting variability and is distinct from the paired image-bootstrap intervals reported in the text.}
    \label{tab:competitor_budget}
\end{table}


\paragraph{More competitors improve fixed-budget transfer.}
We fix 32 records per category and vary the competitor count
\(K\in\{2,4,8\}\), assigning \(32/K\) records to each competitor.
Competitor subsets are nested within each seed, with shared record
permutations and sampling without replacement.
Macro AUC and AP increase monotonically in both eligible cohorts
(\cref{tab:competitor_budget}). Increasing \(K\) from 2 to 8 improves
AUC by \textbf{0.811 percentage points} on Core
(95\% CI \([0.718,0.913]\)) and \textbf{1.379 points} on Extended
(\([1.286,1.468]\)); across-seed AUC SD also decreases.
The Extended comparison retains a positive Bonferroni-adjusted 97.5\% CI
of \([1.274,1.482]\) points. L9 reproduces the mean improvement.
Fixing the per-category counts across three target-area levels to
\((8,8,16)\) preserves the Extended gain on 54 eligible categories
(+1.681 points); no Core category meets this control's support requirements.
Full-record estimates remain stronger on the same cohorts, so these
results support distributing a limited budget across more competitors.

\paragraph{Equal pair weights have a smaller, inconsistent effect.}
Using all records, we compare equal record weights (A0) with equal
competitor weights (A1), preserving raw pair-mean magnitudes and normalizing
only the final direction. A1 changes AUC by \(-0.298\) points on Core
and \(+0.139\) on Extended. Both 95\% CIs lie within the protocol's
\(\pm0.5\)-point practical-effect band. These results do not support a
universal switch to A1; A0 remains the reference estimator.

\paragraph{Family effects depend on the eligible cohort.}
With two competitors and 32 records, cross-family selection exceeds
same-family selection on the 21-category Extended cohort by 6.968 AUC
points, or 5.997 points after area control. The controlled Core difference
is only 0.170 points (95\% CI \([-0.467,0.787]\); three categories).
The Extended cohort is animal-heavy, and changing competitors also changes
donor and background records. This secondary comparison supports a
selection policy within the observed cache, without isolating a causal
effect of semantic distance. The supplementary material gives exact
cohorts, weighting definitions, adjusted intervals, and all controls.
